% Сокращения и обозначения (\listofabbreviations подтягивает их по \newacronym).
% Порядок — по алфавиту (латиница, затем кириллица), как требует ГОСТ 7.32.
\newacronym{aop}{AOP}{Appellation d'Origine Prot\'eg\'ee, защищённое наименование места происхождения (фр., эквивалент PDO)}
\newacronym{api}{API}{application programming interface, программный интерфейс приложения}
\newacronym{bic}{BIC}{Bayesian Information Criterion, байесовский информационный критерий выбора структуры сети}
\newacronym{ci}{CI}{confidence interval, доверительный интервал (95\,\% если не указано иное)}
\newacronym{e2e}{E2E}{end-to-end (сквозной); сквозная точность системы с учётом маршрутизации и отказа от ответа (отказ засчитывается как ошибка); производственно-реалистичная метрика, см. термин «Сквозная точность»}
\newacronym{ece}{ECE}{Expected Calibration Error, ожидаемая ошибка калибровки вероятностей классификатора}
\newacronym{json}{JSON}{JavaScript Object Notation, текстовый формат обмена данными}
\newacronym{lightgbm}{LightGBM}{Light Gradient Boosting Machine, библиотека градиентного бустинга деревьев решений}
\newacronym{llm}{LLM}{large language model, большая языковая модель}
\newacronym{ml}{ML}{machine learning, машинное обучение}
\newacronym{mlp}{MLP}{multi-layer perceptron, многослойный персептрон (полносвязная нейронная сеть)}
\newacronym{mpnet}{MPNet}{Masked and Permuted Pre-training, архитектура трансформера, использованная в paraphrase-multilingual-mpnet-base-v2 для 768-мерных векторных представлений}
\newacronym{odbl}{ODbL}{Open Database License v1.0, лицензия с наследованием (share-alike) для открытых баз данных; применяется к Open Food Facts}
\newacronym{off}{OFF}{Open Food Facts, открытый краудсорсинговый каталог продовольственных товаров под лицензией ODbL}
\newacronym{ood}{OOD}{out-of-distribution, выход за пределы обучающего распределения}
\newacronym{pdo}{PDO}{Protected Designation of Origin, защищённое наименование места происхождения}
\newacronym{pim}{PIM}{product information management, управление информацией о товарах}
\newacronym{rest}{REST}{representational state transfer, архитектурный стиль программного интерфейса}
\newacronym{sbert}{SBERT}{Sentence-BERT, многоязычная модель векторного представления предложений}
\newacronym{sku}{SKU}{stock keeping unit, товарная учётная позиция (артикул)}
\newacronym{tfidf}{TF-IDF}{term frequency $\times$ inverse document frequency, частота терма, взвешенная обратной частотой документа (мера значимости слова в корпусе)}
\newacronym{xgb}{XGBoost}{Extreme Gradient Boosting, библиотека градиентного бустинга деревьев решений}
\newacronym{tz}{ТЗ}{техническое задание}
